\subsection{Księga pierwsza -- postulaty}
\setcounter{counter_euclid}{0}

% Postulat 1
\begin{euclid_postulate}
    Można poprowadzić prostą od któregokolwiek punktu do któregokolwiek punktu.
\end{euclid_postulate}

% Postulat 2.
\begin{euclid_postulate}
    Ograniczoną prostą można przedłużyć nieskończenie.
\end{euclid_postulate}

% Postulat 3.
\begin{euclid_postulate}
    Można zakreślić okrąg z któregokolwiek punktu jako środka dowolną odległością.
\end{euclid_postulate}

% Postulat 4.
\begin{euclid_postulate}
    Wszystkie kąty proste są między sobą równe.
\end{euclid_postulate}

% Postulat 5.
\begin{euclid_postulate}
    Jeżeli prosta przecinająca dwie proste tworzy z nimi kąty jednostronnie wewnętrzne o sumie mniejszej niż dwa kąty proste, to te dwie proste przedłużone nieskończenie przecinają się po tej stronie, po której znajdują się kąty o sumie mniejszej od dwóch kątów prostych.
\end{euclid_postulate}